\usepackage{cleveref} 
% % ******************************** added by Koshiba **********************************
% % ************************************************************************************

\usepackage{threeparttable}
\usepackage{amssymb}
\usepackage{relsize}
\usepackage{amsthm}
\usepackage{here}
\usepackage{algorithm,algorithmic}

% add pages to references
\usepackage[hyperpageref]{backref}
% Setup to show (pages 4 and 9) sort of thing in the bibliography - DD
%\def\foo{\hspace{\fill}\mbox{}\linebreak[0]\hspace*{\fill}}
%\def\foo{\parbox{3cm}{\hfill}
%\def\foo{\parbox{3cm}{\hfill}
%\newcommand\foo[1]{{\raggedleft{\hfill{\mbox{\hfill{#1}}}}}}
\newcommand{\comfyfill}[1]{% = Thorsten Donig's \signed
  \unskip\hspace*{0.1em plus 1fill}
  \nolinebreak[3]%
  \hspace*{\fill}\mbox{#1}
  \parfillskip0pt\par
}
\newcommand\foo[1]{{\comfyfill{\mbox{#1}}}}
%\newcommand\foo[1]{{\mbox{#1}}}
\renewcommand*{\backref}[1]{}
\renewcommand*{\backrefalt}[4]{%
\ifcase #1 %
%
\or
\foo{(page #2)}%
\else
\foo{(pages #2)}%
\fi
}

% add colors for linsk 
\usepackage{xcolor}
\hypersetup{
    colorlinks=true,
    citecolor=[rgb]{0, 0, 0.6},
    linkcolor=red,
    urlcolor=[rgb]{0, 0, 0.6},
}

% Proof stuff
\theoremstyle{plain}
\newtheorem{definition}[]{Definition}
\newtheorem{theorem}{Theorem}[section]
\newtheorem{lemma}[theorem]{Lemma}
\newtheorem{prop}[theorem]{Proposition}
\newtheorem{proposition}{Proposition}
\newtheorem*{cor}{Corollary}

% abbreviated words
\newcommand{\ds}{$D_s$\ }
\newcommand{\cm}{m$^3$}
\newcommand{\fr}{\mathrm{Fr}}
\newcommand{\ips}[1]{I_{ps}^{#1}}
\newcommand{\Ker}[2]{\mathrm{Ker}(#1;\,#2)}
\newcommand{\vips}{\boldsymbol{I_{ps}}}


% for math
\newcommand{\pr}{{\mathrm{Pr}}}
\newcommand{\vA}{\mathbf{A}}
\newcommand{\vI}{\mathbf{I}}
\newcommand{\vk}{\mathbf{k}}
\newcommand{\vK}{\mathbf{K}}
\newcommand{\vLambda}{\mathbf{\Lambda}}
\newcommand{\vmu}{{\boldsymbol{\mu}}}
\newcommand{\vphi}{{\boldsymbol{\phi}}}
\newcommand{\vPhi}{\mathbf{\Phi}}
\newcommand{\vsigma}{{\boldsymbol{\sigma}}}
\newcommand{\vSigma}{\mathbf{\Sigma}}
\newcommand{\vtheta}{{\boldsymbol{\theta}}}
\newcommand{\vTheta}{\mathbf{\Theta}}
\newcommand{\vW}{\mathbf{W}}
\newcommand{\vw}{\mathbf{w}}
\newcommand{\vX}{\mathbf{X}}
\newcommand{\vx}{\mathbf{x}}
\newcommand{\vY}{\mathbf{Y}}
\newcommand{\vy}{\mathbf{y}}
\newcommand{\vzero}{\boldsymbol{0}}
\newcommand{\vZ}{\mathbf{Z}}
\newcommand{\vz}{\mathbf{z}}

\newcommand{\Normal}{\mathcal{N}}
\newcommand{\tra}{^{\mathsf{T}}}
\newcommand{\inv}{^{{\mathsmaller{-1}}}}
\newcommand{\expect}{\mathbb{E}}
\newcommand{\expectargs}[2]{\mathbb{E}_{#1} \left[ {#2} \right]}
\newcommand{\var}{\mathbb{V}}
\newcommand{\varianceargs}[2]{\mathbb{V}_{#1} \left[ {#2} \right]}
\newcommand{\gpcal}{\mathcal{GP}}
%%%% others
% \newcommand{\figwid}{0.47}
\newcommand{\figwid}{7.2cm}
\newcommand{\cs}[1]{\texttt{#1}}

% for tabular photos
\newcommand{\ttwop}[1]{\includegraphics[width=7.2 cm]{#1}}
\newcommand{\tthrp}[1]{\includegraphics[width=4.8 cm]{#1}}
\newcommand{\tfoup}[1]{\includegraphics[width=\figwid]{#1}}

% % FOR TABLES
% for multi-row table elements
\newcommand{\onerow}[1]{\begin{tabular}{l}#1\end{tabular}}
\newcommand{\tworow}[2]{\begin{tabular}{l}#1\\#2\end{tabular}}
\newcommand{\trirow}[3]{\begin{tabular}{l}#1\\#2\\#3\end{tabular}}
% for footnotes
\newcommand{\f}{\tnote{*}}
\newcommand{\ff}{\tnote{**}}
\newcommand{\fff}{\tnote{***}}